% =============================================================================
%  Chapter 1 — Simplex Finite Elements
% =============================================================================

\chapter{Simplex Finite Elements}
\label{ch:simplex}

This chapter describes the construction of Lagrangian finite elements on
simplicial cells (triangles in 2D, tetrahedra in 3D) as implemented in
\texttt{fall\_n}.  The development follows the classical Ciarlet
definition~\cite{Ciarlet1972} specialised to Lagrange interpolation on the
reference simplex.


% ─────────────────────────────────────────────────────────────────────────
\section{Reference Simplex}
\label{sec:refsimplex}

\begin{definition}[Reference simplex]
The \emph{reference simplex} in $\R^D$ is the convex hull
\begin{equation}
  \calS_D \;=\; \Bigl\{
    \bxi = (\xi_1, \dots, \xi_D) \in \R^D
    \;\Big|\;
    \xi_i \ge 0,\;
    \sum_{i=1}^{D} \xi_i \le 1
  \Bigr\},
  \label{eq:ref_simplex}
\end{equation}
with \emph{vertices}
\[
  \mathbf{v}_0 = \mathbf{0}, \qquad
  \mathbf{v}_i = \mathbf{e}_i \quad (i = 1, \dots, D),
\]
where $\mathbf{e}_i$ is the $i$-th standard basis vector.
\end{definition}

The \emph{barycentric coordinates} $(\lambda_0, \lambda_1, \dots, \lambda_D)$
are defined as
\begin{equation}
  \lambda_0 = 1 - \sum_{i=1}^{D} \xi_i, \qquad
  \lambda_i = \xi_i \quad (i = 1, \dots, D).
  \label{eq:bary}
\end{equation}
They form a partition of unity, $\sum_{i=0}^{D} \lambda_i = 1$, and satisfy
$\lambda_i \ge 0$ on $\calS_D$.

The volume of the reference simplex is
\begin{equation}
  |\calS_D| = \frac{1}{D!}\,.
  \label{eq:simplex_volume}
\end{equation}


% ─────────────────────────────────────────────────────────────────────────
\section{Lagrange Shape Functions}
\label{sec:shape_functions}

\subsection{Linear Elements (Order 1)}

The linear simplex has $D+1$ nodes (the vertices).  The shape functions are
simply the barycentric coordinates:
\begin{equation}
  N_i(\bxi) = \lambda_i(\bxi), \qquad i = 0, 1, \dots, D.
  \label{eq:shape_linear}
\end{equation}

Since $\lambda_i \ge 0$ on $\calS_D$ and $\lambda_i(\mathbf{v}_j) =
\delta_{ij}$, these are the unique linear polynomials that interpolate the
nodal values.

\paragraph{Concrete instances in \texttt{fall\_n}:}
\begin{itemize}[nosep]
  \item TRI3 ($D=2$, 3 nodes): \texttt{SimplexElement<3,2,1>} (triangle in 3D
        ambient space, used for boundary faces)
  \item TET4 ($D=3$, 4 nodes): \texttt{SimplexElement<3,3,1>}
\end{itemize}


\subsection{Quadratic Elements (Order 2)}
\label{sec:quadratic_simplex}

The quadratic simplex has
\begin{equation}
  n_{\text{nodes}} = \binom{D+2}{2} = \frac{(D+1)(D+2)}{2}
  \label{eq:numnodes_quadratic}
\end{equation}
nodes: the $D+1$ vertices plus $\binom{D+1}{2}$ edge midpoints.

\paragraph{Vertex shape functions.}
For vertex $i$ ($i=0,\dots,D$):
\begin{equation}
  \boxed{N_i(\bxi) = \lambda_i\,(2\lambda_i - 1)}
  \label{eq:shape_vertex_quad}
\end{equation}

\begin{remark}[Sign of vertex shape functions]
\label{rem:vertex_negative}
The vertex shape function $N_i = \lambda_i(2\lambda_i - 1)$ is \emph{negative}
whenever $0 < \lambda_i < \tfrac{1}{2}$.  This region is non-empty in the
interior of $\calS_D$.  In particular, interior quadrature points that are
``far'' from vertex $i$ will have $\lambda_i \ll \tfrac{1}{2}$, making
$N_i < 0$ at those points.

This fact is benign for stiffness matrix assembly (the bilinear form involves
gradients, not function values).  However, it becomes critical for lumped $L^2$
projection, as discussed in~\cref{ch:negative_weights}.
\end{remark}

\paragraph{Edge midpoint shape functions.}
For the midpoint of edge $(i, j)$ with $0 \le i < j \le D$:
\begin{equation}
  \boxed{N_{ij}(\bxi) = 4\,\lambda_i\,\lambda_j}
  \label{eq:shape_edge_quad}
\end{equation}

These are always non-negative on $\calS_D$, since both $\lambda_i, \lambda_j
\ge 0$.

\paragraph{Concrete instances:}
\begin{itemize}[nosep]
  \item TRI6 ($D=2$, 6 nodes): \texttt{SimplexElement<3,2,2>}
  \item TET10 ($D=3$, 10 nodes): \texttt{SimplexElement<3,3,2>}
\end{itemize}


% ─────────────────────────────────────────────────────────────────────────
\section{Shape Function Derivatives}
\label{sec:derivatives}

Since $\lambda_0 = 1 - \xi_1 - \cdots - \xi_D$ and $\lambda_i = \xi_i$ for
$i \ge 1$, the barycentric coordinate gradients (with respect to $\bxi$) are:
\begin{equation}
  \frac{\partial \lambda_0}{\partial \xi_k} = -1, \qquad
  \frac{\partial \lambda_i}{\partial \xi_k} = \delta_{ik}
  \quad (i = 1, \dots, D).
  \label{eq:bary_grad}
\end{equation}

For the quadratic vertex shape function $N_i = \lambda_i(2\lambda_i - 1)$:
\begin{equation}
  \frac{\partial N_i}{\partial \xi_k}
  = (4\lambda_i - 1)\,\frac{\partial \lambda_i}{\partial \xi_k}\,.
  \label{eq:dNi_vertex}
\end{equation}

For the edge midpoint $N_{ij} = 4\lambda_i\lambda_j$:
\begin{equation}
  \frac{\partial N_{ij}}{\partial \xi_k}
  = 4\Bigl(
    \lambda_j\,\frac{\partial \lambda_i}{\partial \xi_k}
  + \lambda_i\,\frac{\partial \lambda_j}{\partial \xi_k}
  \Bigr).
  \label{eq:dNij_edge}
\end{equation}


% ─────────────────────────────────────────────────────────────────────────
\section{Isoparametric Mapping}
\label{sec:isoparametric}

The isoparametric map from the reference simplex $\calS_D$ to a physical
element $\Omega^e \subset \R^d$ (with $d \ge D$) is
\begin{equation}
  \mathbf{x}(\bxi)
  = \sum_{I=1}^{n_{\text{nodes}}} N_I(\bxi)\,\mathbf{x}_I,
  \label{eq:isoparametric_map}
\end{equation}
where $\mathbf{x}_I \in \R^d$ are the physical nodal coordinates.  The
\emph{Jacobian matrix} of the mapping is
\begin{equation}
  J_{ij} = \frac{\partial x_i}{\partial \xi_j}
         = \sum_{I=1}^{n_{\text{nodes}}}
           \frac{\partial N_I}{\partial \xi_j}\, x_{I,i}\,,
  \qquad
  \mathbf{J} \in \R^{d \times D}.
  \label{eq:jacobian}
\end{equation}

When $d = D$ (volume elements), $\mathbf{J}$ is square and the differential
volume element is $\dd\Omega = |\det\mathbf{J}|\,\dd\bxi$.

When $d > D$ (surface elements, e.g.\ TRI6 boundary faces in 3D), the surface
measure is computed from the metric tensor:
\begin{equation}
  \dd S = \sqrt{\det(\mathbf{J}^T \mathbf{J})}\;\dd\bxi.
  \label{eq:surface_measure}
\end{equation}

\paragraph{Gradient transformation.}
The gradient of a scalar field in physical space is obtained from the
reference-space gradient via
\begin{equation}
  \frac{\partial N_I}{\partial x_j}
  = \sum_{k=1}^{D} \frac{\partial N_I}{\partial \xi_k}\,
    (J^{-1})_{kj}\,.
  \label{eq:grad_transform}
\end{equation}


% ─────────────────────────────────────────────────────────────────────────
\section{Subentity Topology: Faces of the Simplex}
\label{sec:faces}

A $D$-simplex has $D+1$ faces (codimension-$1$ sub-simplices).  The $f$-th face
($f = 0, \dots, D$) is the sub-simplex \emph{opposite} vertex $f$, spanned by
all vertices \emph{except} $\mathbf{v}_f$.

\subsection{Face Local Node Ordering (TET10)}

For a quadratic tetrahedron ($D=3$, TET10, 10~nodes), each of the 4~triangular
faces is a TRI6 sub-element.  The face-to-global node mapping implemented in
\texttt{SimplexCell<3,2>} is:

\begin{center}
\begin{tabular}{ccc}
  \toprule
  Face $f$ & Opposite vertex & Global nodes \\
  \midrule
  0 & $\mathbf{v}_0$ & (1, 2, 3, 5, 8, 7) \\
  1 & $\mathbf{v}_1$ & (0, 3, 2, 6, 8, 5) \\
  2 & $\mathbf{v}_2$ & (0, 1, 3, 4, 7, 6) \\
  3 & $\mathbf{v}_3$ & (0, 2, 1, 5, 4, 6) \\
  \bottomrule
\end{tabular}
\end{center}

The outward normal direction follows from the right-hand rule applied to the
first three (vertex) nodes of each face.  This topology is critical for
consistent surface traction application via Gauss quadrature on boundary faces.


% ─────────────────────────────────────────────────────────────────────────
\section{Gmsh Node Ordering for TET10}
\label{sec:gmsh_reorder}

Gmsh uses a different midpoint ordering than
\texttt{fall\_n}~\cite{Geuzaine2009}.  The Gmsh TET10 local ordering is:

\begin{center}
\begin{tabular}{ccccccccccc}
  \toprule
  Position & 0 & 1 & 2 & 3 & 4 & 5 & 6 & 7 & 8 & 9 \\
  \midrule
  Gmsh     & $v_0$ & $v_1$ & $v_2$ & $v_3$ & $m_{01}$ & $m_{12}$ & $m_{02}$ & $m_{03}$ & $m_{13}$ & $m_{23}$ \\
  \texttt{fall\_n} & $v_0$ & $v_1$ & $v_2$ & $v_3$ & $m_{01}$ & $m_{02}$ & $m_{03}$ & $m_{12}$ & $m_{23}$ & $m_{13}$ \\
  \bottomrule
\end{tabular}
\end{center}

The reordering applied in \texttt{GmshDomainBuilder.hh} maps Gmsh indices to
\texttt{fall\_n} indices via the permutation:
\begin{equation}
  (0, 1, 2, 3, 4, 6, 7, 5, 9, 8).
  \label{eq:gmsh_perm}
\end{equation}
This is implemented as:
\begin{lstlisting}
index_ordering(node_tags, 0,1,2,3,4,6,7,5,9,8).data()
\end{lstlisting}


% ─────────────────────────────────────────────────────────────────────────
\section{Implementation in \texttt{fall\_n}}
\label{sec:simplex_impl}

The simplex element hierarchy consists of three layers:

\begin{enumerate}[leftmargin=2em]
  \item \textbf{\texttt{SimplexCell<D, Order>}} (\texttt{SimplexCell.hh}) ---
        Purely geometric: reference coordinates, shape functions and their
        derivatives, face topology.  All methods are \texttt{consteval} or
        \texttt{constexpr}, enabling compile-time evaluation.

  \item \textbf{\texttt{SimplexElement<Dim, TopDim, Order>}}
        (\texttt{SimplexElement.hh}) --- Combines a \texttt{SimplexCell} with
        an integration strategy.  Template parameters:
        \begin{itemize}[nosep]
          \item \texttt{Dim}: ambient space dimension (e.g., 3 for 3D problems);
          \item \texttt{TopDim}: topological dimension ($D$);
          \item \texttt{Order}: Lagrange interpolation order.
        \end{itemize}

  \item \textbf{\texttt{ElementGeometry<ElemType, Integrator>}}
        (\texttt{ElementGeometry.hh}) --- Manages the Jacobian, physical
        gradients, and integration; parameterised by element type and
        quadrature rule.  The integrator is a duck-typed object exposing
        \texttt{num\_integration\_points}, \texttt{reference\_integration\_point(i)},
        and \texttt{weight(i)}.
\end{enumerate}

This layered design allows the quadrature rule to be swapped independently of
the element geometry, which was crucial for diagnosing and resolving the
negative-weight problem (see~\cref{ch:negative_weights,ch:quadrature}).
